\begin{center}
{\large \bf \TITLE}
\end{center}

\begin{center}
    {\bf Synergistic Activities}
\end{center}

\begin{enumerate}
    \itemsep=-1pt
    \item {\bf Policy Leadership and Public Engagement:} Feamster leads the University of Chicago's AI and Tech Policy Pillar, coordinating efforts across campus to address critical questions at the intersection of artificial intelligence and public policy. He serves on the Advisory Board of the University of Chicago Forum on Free Expression, contributing expertise on online speech, content moderation, and digital communication platforms. Feamster works to ensure that technical research results inform policy through participation and leadership in national policy organizations (e.g., Broadband Internet Technology Advisory Group, where he has co-authored multiple technical reports). He testified before the General Accountability Office (GAO) on Internet of Things security and has provided extensive public comments to the FCC on privacy and net neutrality rulemakings, with his comments cited more than 30 times in FCC rulemakings.

    \item {\bf Research Leadership and Translation to Practice:} Feamster serves as Director of Research and Tech Policy at the University of Chicago Data Science Institute and directs the Internet Equity Initiative, which has secured more than \$2M in research funding. Through these roles, he has led the development and deployment of the Netrics open-source measurement platform, now used by approximately ten research groups and several state governments for broadband measurement. The initiative's work has directly informed state and federal broadband policy, including the creation of the BEAD Challenge platform (\url{https://beadchallenge.org/}), which enables citizens to challenge FCC broadband maps and has thousands of daily visitors. Feamster has also transitioned research to practice through NetMicroscope, Inc., securing \$500K in customer revenue and filing patent applications.

    \item {\bf Conference Leadership:} Feamster founded the USENIX Workshop on Free and Open Communication on the Internet (FOCI) and continues to lead its steering committee. The workshop is now in its 13th year and serves as the premier venue for research on Internet censorship and information control. He has served as Program Committee Co-Chair for major conferences including ACM SIGCOMM Internet Measurement Conference (IMC), USENIX Symposium on Networked Systems Design and Implementation (NSDI), and Internet Engineering Task Force Applied Networking Research Workshop (ANRW). He serves on the steering committees for NSDI, ANRW, and the ACM SIGCOMM Workshop on Networking Meets AI/ML.

    \item {\bf Education and Curriculum Development:} Feamster serves as Faculty Director for the University of Chicago Cybersecurity Professional Education Program, providing advanced training in cybersecurity to working professionals. He teaches courses on Internet Censorship and Online Speech to the University of Chicago Masters of Liberal Arts Program and through the Parrhesia program on free expression. He teaches a Machine Learning for Systems course at the University of Chicago's International Institute for Research in Paris (IIRP) as part of the study abroad program, focusing on generative AI applications for networking. Feamster has demonstrated exceptional commitment to broad educational dissemination: his YouTube courses on Computer Networking, Software Defined Networking, and Internet Censorship and Online Speech have collectively garnered approximately one million views. He developed the Coursera Software Defined Networking (SDN) Course, which had 50,000 enrollments in its first offering (2013) and has had more than 200,000 learners in total. He is currently writing two textbooks: {\em Machine Learning for Computer Systems} (with Bruce Davie, Systems Approach LLC) and {\em Internet Censorship and Online Speech}.

    \item {\bf Broadening Participation in Computing:} Feamster teaches an annual ``Applications of Machine Learning to Networking'' summer course through the University of Chicago's Upward Bound program to underrepresented students across the Chicago Public School system (77\% economically disadvantaged, 83\% underrepresented minority). This program includes an annual series of 10 seminars and 6 weekend workshops. He has mentored two high school students from Walter Payton College Prep on Advanced Placement Research projects, with one becoming a first-generation college student at the University of Chicago. Feamster has successfully mentored numerous Ph.D. students who have gone on to faculty positions at top institutions, including women and underrepresented minorities.
\end{enumerate}
